\section{LLM-as-a-judge prompt}
\label{app:prompt-llm-judge}

\subsection{Meaning-representation pipeline prompts}
\label{app:pipeline-judge-prompts}

The meaning-representation pipeline evaluation uses four independent criteria,
split across two transitions. The data-to-text transition uses summary and
faithfulness; the text-to-triples transition uses additions and omissions. Each
criterion is scored with a separate prompt on an integer scale from 1 (worst)
to 5 (best). The judge is instructed to assess only the relationship between
the supplied source and output, without using external knowledge:

\begin{description}
  \item[Data to text] The source is the original structured data instance and
  the output is the generated natural-language reference text.
  \item[Text to triples] The source is the generated reference text and the
  output is the generated semantic triples.
\end{description}

The task-specific part of the system prompt is instantiated as follows:

\begin{lstlisting}
You are a strict judge evaluating a {conversion} conversion.
You will see: (1) {source_label}, and (2) {output_label}.

{criterion_guideline}

Use only the provided source and output. Assign one integer score from 1 to 5.
Return ONLY JSON with this exact schema:
{"score": <integer 1-5>, "reason": "<one short sentence>"}
\end{lstlisting}

For the data-to-text transition, \texttt{conversion} is ``data-to-text'',
\texttt{source\_label} is ``the original structured data instance'', and
\texttt{output\_label} is ``the natural-language text generated from it''.
For the text-to-triples transition, the corresponding values are
``text-to-triples'', ``the natural-language reference text'', and ``the
semantic triples generated from it''. The current criterion guidelines are:

\subsubsection{Summary}

\begin{lstlisting}
Criterion: Summary
Evaluate whether the natural-language text is a concise, coherent summary of
the original structured data instance. It should avoid repeating the same
concept with different words, be significantly shorter than detailed or
repetitive structured data when possible, preserve the core message, and
function as a summary rather than a rewritten essay or a disjointed list of
facts.

Score 5: The natural-language text is a clear and coherent summary. It
preserves the original structured data's core message without redundant
restatement and does not read like an essay or a disconnected list.
Score 4: The natural-language text is a good summary and preserves the core
message, but has minor redundancy, slight unnecessary length, or a small
organizational weakness.
Score 3: The natural-language text is recognizably a summary, but is noticeably
repetitive, too verbose, list-like, or weakly organized; its main message
remains usable.
Score 2: The natural-language text is a poor summary: it substantially repeats
or rewrites the original structured data, is very verbose or disjointed, and
represents the core message only partly.
Score 1: The natural-language text does not function as a summary. It is mostly
irrelevant, repetitive, disjointed, or fails to communicate the original
structured data's core message.
\end{lstlisting}

\subsubsection{Additions}

\begin{lstlisting}
Criterion: Additions
Evaluate whether the semantic triples introduce facts, relationships, values,
qualifiers, or conclusions that are not supported by the natural-language
reference text. Penalize unsupported additions, hallucinations, contradictions,
and material distortions. Do not penalize information that is merely omitted;
evaluate omissions under Omissions.

Score 5: The semantic triples contain no unsupported additions. Every
represented fact, relationship, value, qualifier, and conclusion is supported
by the natural-language reference text, with no hallucination or material
distortion.
Score 4: The semantic triples contain at most one or two minor unsupported
additions or imprecisions, but no material fact is fabricated and the text's
meaning is preserved.
Score 3: The semantic triples contain several minor unsupported additions or
one material unsupported fact or distortion, but most represented information
is grounded in the text.
Score 2: The semantic triples contain multiple material unsupported facts,
relationships, values, or conclusions, or substantially misrepresent the
reference text.
Score 1: The semantic triples are mostly unsupported, contradictory, or
materially different from the natural-language reference text.
\end{lstlisting}

\subsubsection{Faithfulness}

\begin{lstlisting}
Criterion: Faithfulness
Evaluate whether every claim in the natural-language text is directly supported
by the original structured data instance and whether entities, relationships,
values, qualifiers, and their meaning are represented accurately. Penalize
hallucinations, fabrications, contradictions, and misinterpretations. Do not
penalize information that is merely omitted.

Score 5: Every claim is grounded in the original structured data instance, and
all represented facts preserve the source's meaning. There are no
hallucinations, contradictions, or material distortions.
Score 4: The natural-language text is fully usable and essentially faithful,
with at most a minor imprecision that does not change the meaning and no
material unsupported claim.
Score 3: Most claims are grounded, but there is one material error or
distortion, or several minor unsupported or inaccurate claims.
Score 2: There are multiple material errors, hallucinations, contradictions, or
a substantial misinterpretation of the original structured data.
Score 1: The natural-language text is mostly ungrounded or contradicts the
original structured data, so it cannot be considered a reliable representation.
\end{lstlisting}

\subsubsection{Omissions}

\begin{lstlisting}
Criterion: Omissions
Evaluate what the semantic triples leave out from the output text. Lower the
score when omissions remove critical context, change the meaning, hide a major
conclusion or constraint, or create a misleading imbalance or bias.

Score 5: Only redundant material is omitted. No omission changes the meaning,
hides important context, or introduces a meaningful bias.
Score 4: One or a few minor omissions exist, but the meaning and conclusions
remain intact and the selection of included information is not meaningfully
biased.
Score 3: An important but non-central detail, or enough supporting context to
create a mild imbalance, is omitted; the overall message remains
understandable.
Score 2: A critical context element, major conclusion, important constraint, or
representative point is omitted, changing the meaning or creating clear bias.
Score 1: Extensive omissions distort the core message or make the output
triples misleading or one-sided; most important source information is absent.
\end{lstlisting}

\subsection{Surface-realization judge prompts}
\label{app:surface-judge-prompts}

For every non-empty generated realization, the evaluator sends the input
triples and the generated text to the configured LLM judge. The reference texts
are not included in this request. The prompt used by the evaluator is shown
below; \texttt{source} is the formatted list of triples and \texttt{target} is
the generated text.

\begin{lstlisting}
###Instruction###
Please act as an impartial and helpful evaluator for natural language generation (NLG), and the audience is an expert in the field.
Your task is to evaluate the quality of text similarity strictly based on the given evaluation criterion.
Begin the evaluation by providing your analysis concisely and accurately, and then on the next line, start with "Rating:" followed by your rating on a Likert scale from 1 to 5 (higher means better).
You MUST keep to the strict boundaries of the evaluation criterion and focus solely on the issues and errors involved; otherwise, you will be penalized.
Make sure you read and understand these instructions, as well as the following evaluation criterion and example content, carefully.

###Evaluation Criterion###
Accuaracy and number of sentences: The generated text must accurately reference the triples. There cannot be any extra information that was not present in the triples. If possible there must be just one complex sentence instead of multiple sentences.

###Data###
The triples in the form (subject, predicate, object):
{source}

The generated text:
{target}

###Your Evaluation###
\end{lstlisting}

The response is parsed as a review followed by a rating from 1 to 5. The
ratings are averaged and divided by 5 before they are used as the evolutionary
fitness.

\subsection{Grammaticality prompt}

The following criterion is used for the binary grammaticality score. The
judge returns 1 for a grammatically correct realization and 0 otherwise.

\begin{lstlisting}
Grammatical correctness: You should assess the grammatical correctness of the
resulting text. Do not take any other factors into account. Do not make
assumptions or consider external knowledge not present in the provided context.
Identify only errors relating to the grammaticality of the text. Do not
consider aspects such as fluency, omissions or hallucinations.

Return 1 if the text is grammatically correct and 0 if it is not. The rating
must be the integer 0 or 1.
\end{lstlisting}

\subsection{Omissions prompt}

The omissions score measures whether any input triples were not verbalized. A
score of 0 indicates that no omission was detected, while 1 indicates that at
least one omission was detected.

\begin{lstlisting}
Omissions: You should assess the omissions in the resulting text; in other
words, check whether any of the input triples were not verbalised. You can
perform the task by iterating over the input triples and checking if each is
present in the output. Do not take any other factors into account. Do not make
assumptions or consider external knowledge not present in the provided context.
Identify only omissions. Do not consider grammaticality, fluency or the
addition of new facts.

Return 0 if there are no omissions and 1 if there are. The rating must be the
integer 0 or 1.
\end{lstlisting}

\subsection{Additions prompt}

The additions score measures whether the realization introduces facts that are
not supported by the input triples. A score of 0 indicates that no addition was
detected, while 1 indicates that at least one addition was detected.

\begin{lstlisting}
Additions: You should assess the addition of new facts in the resulting text
which were not present in the input triples. Carefully read the text and check
whether the facts mentioned are present in the input triples. Do not take any
other factors into account. Do not make assumptions or consider external
knowledge not present in the provided context. Identify only additions. Do not
consider grammaticality, fluency or omissions.

Return 0 if there are no additions and 1 if there are. The rating must be the
integer 0 or 1.
\end{lstlisting}

\subsection{Zero-shot prompt}

The zero-shot prompt is used to convert the input triples into a natural language realization without any additional context. 
The prompt is shown below;

\begin{lstlisting}
You are an expert data-to-text generator. Verbalize the given RDF triples 
into a single fluent, grammatically correct English sentence that conveys 
every fact and adds no extra information.

Convert these RDF triples (subject | predicate | object) into one
natural-language sentence.

Triples:
{triples}

Output only the sentence.
\end{lstlisting}
